% 福州大学本科生毕业设计（论文）开题报告
% 填写入口：
% 1. 先修改“基本信息”中的姓名、学号、专业、题目。
% 2. 再按 1 到 7 项替换正文提示文本，写入你的开题报告内容。
% 3. “指导教师评语”和下方签字日期通常由教师填写。
% 编译方式：latexmk -xelatex -shell-escape -interaction=nonstopmode -halt-on-error main.tex

\documentclass[12pt,a4paper]{ctexart}
\usepackage{geometry}
\usepackage{longtable}
\usepackage{array}
\usepackage{adjustbox}
\usepackage{multirow}
\usepackage{graphicx}
\usepackage{etoolbox}
\usepackage{amssymb}

\setlength{\parindent}{0pt}
\setlength{\tabcolsep}{0pt}
\setlength{\arrayrulewidth}{0.4pt}
\renewcommand{\arraystretch}{1.0}

\newcommand{\FZUCenteredUnderline}[2]{%
  \underline{\makebox[#1][c]{#2}}%
}

\newcommand{\FZUCellBox}[2]{%
  \begin{minipage}[t][#1][t]{\linewidth}%
    \setlength{\parindent}{0pt}%
    #2%
  \end{minipage}%
}

\newsavebox{\FZUTextFitBox}
\newcommand{\FZUFitTextBlock}[2]{%
  \sbox{\FZUTextFitBox}{%
    \begin{minipage}[t]{\linewidth}%
      \setlength{\parindent}{0pt}%
      \raggedright
      #2%
    \end{minipage}%
  }%
  \ifdim\dimexpr\ht\FZUTextFitBox+\dp\FZUTextFitBox\relax>#1\relax
    \resizebox{!}{#1}{\usebox{\FZUTextFitBox}}%
  \else
    \usebox{\FZUTextFitBox}%
  \fi
}

\newcommand{\FZUCheckboxBox}[1]{%
  \makebox[1.05em][c]{%
    \raisebox{0.05em}{%
      \fboxsep=0pt%
      \framebox[0.86em][c]{\rule{0pt}{0.84em}#1}%
    }%
  }%
}

\newcommand{\FZUCheckboxEmpty}{\FZUCheckboxBox{}}
\newcommand{\FZUCheckboxChecked}{\FZUCheckboxBox{\smash{\raisebox{0.06em}{\scriptsize$\checkmark$}}}}


\geometry{top=2.1cm,bottom=2.05cm,left=2.35cm,right=2.35cm}

% ===== 基本信息 =====
\newcommand{\姓名}{}
\newcommand{\学号}{}
\newcommand{\专业}{}
\newcommand{\题目}{}
\newcommand{\指导教师评语}{此处由指导教师填写对开题报告的评语。}
\newcommand{\签字年}{}
\newcommand{\签字月}{}
\newcommand{\签字日}{}
% ====================

\newlength{\ProposalWidth}
\newlength{\ProposalInnerWidth}
\setlength{\ProposalWidth}{14.00cm}
\setlength{\ProposalInnerWidth}{13.54cm}
\setlength{\LTleft}{0pt}
\setlength{\LTright}{0pt}
\setlength{\LTpre}{0pt}
\setlength{\LTpost}{0pt}

\newcommand{\ProposalTitleRow}[1]{%
  \hspace*{0.23cm}%
  \begin{minipage}[t]{\ProposalInnerWidth}%
    \setlength{\parindent}{0pt}%
    \raggedright
    \vspace{0.12cm}%
    {\zihao{5}\bfseries #1}%
    \vspace{0.05cm}%
  \end{minipage}\tabularnewline
}

\newcommand{\ProposalTextRow}[1]{%
  \hspace*{0.23cm}%
  \begin{minipage}[t]{\ProposalInnerWidth}%
    \setlength{\parindent}{2em}%
    \raggedright
    \vspace{0.06cm}%
    {\zihao{5}#1}%
    \vspace{0.08cm}%
  \end{minipage}\tabularnewline
}

\newcommand{\ProposalNoteRow}[1]{%
  \hspace*{0.23cm}%
  \begin{minipage}[t]{\ProposalInnerWidth}%
    \setlength{\parindent}{0pt}%
    \raggedright
    \vspace{0.10cm}%
    {\zihao{5}\bfseries #1}%
    \vspace{0.12cm}%
  \end{minipage}\tabularnewline
}

\newcommand{\ProposalIntroRow}{%
  \hspace*{0.18cm}%
  \begin{minipage}[t]{13.64cm}%
    \setlength{\parindent}{0pt}%
    \raggedright
    \vspace{0.12cm}%
    {\zihao{5}开题报告应包括下列主要内容：}\par
    \vspace{0.08cm}
    {\zihao{5}\hspace*{1.8em}1．课题来源及研究的目的和意义；2．国内外在该方向的研究现状及分析；3．主要研究内容；4．研究方案及进度安排，预期达到的目标；5．为完成课题已具备和所需的条件；6．预计研究过程中可能遇到的困难和问题，以及解决的措施；7．主要参考文献。}\par
    \vspace{0.14cm}%
  \end{minipage}\tabularnewline
}

\begin{document}

\begin{center}
  \vspace*{0.05cm}
  {\zihao{3}\bfseries 福州大学本科生毕业设计（论文）开题报告}
\end{center}

\vspace{0.18cm}
\noindent
\begin{longtable}{|p{\ProposalWidth}|}
\hline
\begin{tabular}{>{\centering\arraybackslash}p{1.08cm}|p{1.45cm}|>{\centering\arraybackslash}p{1.08cm}|p{1.88cm}|>{\centering\arraybackslash}p{1.18cm}|p{7.33cm}}
\hline
\rule{0pt}{0.92cm}{\zihao{5}\bfseries 姓名} & \zihao{5}\姓名 &
\zihao{5}\bfseries 学号 & \zihao{5}\学号 &
\zihao{5}\bfseries 专业 & \zihao{5}\专业 \\
\hline
\multicolumn{2}{>{\centering\arraybackslash}p{2.53cm}|}{\rule{0pt}{0.92cm}\zihao{5}\bfseries 毕业设计（论文）题目} &
\multicolumn{4}{p{11.47cm}}{\zihao{5}\hspace{0.3em}\题目} \\
\hline
\end{tabular}%
\tabularnewline
\ProposalIntroRow
\hline
\endfirsthead
\hline
\endhead
\hline
\endfoot
\hline
\endlastfoot
\ProposalTitleRow{1．课题来源及研究的目的和意义}
\ProposalTextRow{此处填写课题来源、研究背景、现实意义与学术价值，说明为什么要做这个题目，以及它要解决什么问题。}
\ProposalTitleRow{2．国内外在该方向的研究现状及分析}
\ProposalTextRow{此处梳理国内外相关研究进展，可概括代表性方法、已有成果、当前不足，以及本文准备切入的研究位置。}
\ProposalTitleRow{3．主要研究内容}
\ProposalTextRow{此处概括论文将完成的主要工作，建议分点说明数据准备、方法设计、系统实现、实验验证或其他核心任务。}
\ProposalTitleRow{4．研究方案及进度安排，预期达到的目标}
\ProposalTextRow{此处填写研究方案、技术路线、阶段进度安排和预期目标。建议写清时间节点、阶段成果和最终交付物。}
\ProposalTitleRow{5．为完成课题已具备和所需的条件}
\ProposalTextRow{此处说明目前已经具备的条件，例如文献基础、数据条件、实验环境、开发基础、指导支持，以及后续仍需补充的资源。}
\ProposalTitleRow{6．预计研究过程中可能遇到的困难和问题，以及解决的措施}
\ProposalTextRow{此处填写可能遇到的技术难点、数据问题、进度风险或实验风险，并对应给出拟采取的解决措施。}
\ProposalTitleRow{7．主要参考文献}
\ProposalTextRow{此处填写开题报告参考文献。建议列出与课题直接相关的基础文献、代表性方法文献和应用文献，并按学校要求统一格式。}
\ProposalNoteRow{（不够写可另加页）}
\end{longtable}

\vspace{0.18cm}
\noindent
\begin{tabular}{|p{\ProposalWidth}|}
\hline
\hspace*{0.23cm}\begin{minipage}[t][5.15cm][t]{\ProposalInnerWidth}
\setlength{\parindent}{0pt}%
\raggedright
\vspace{0.12cm}%
{\zihao{5}指导教师评语：}\par
\vspace{0.14cm}
{\zihao{5}\指导教师评语}\par
\vfill
\end{minipage}\tabularnewline
\hline
\end{tabular}

\vspace{0.26cm}
\noindent
\begin{tabular}{>{\centering\arraybackslash}p{2.35cm}p{3.25cm}>{\centering\arraybackslash}p{1.95cm}p{4.75cm}}
\zihao{5}\bfseries 指导教师签字 & & \zihao{5}\bfseries 签字时间 &
\zihao{5}\签字年\hspace{0.8em}年\hspace{0.8em}\签字月\hspace{0.8em}月\hspace{0.8em}\签字日\hspace{0.8em}日
\end{tabular}

\vspace{0.16cm}
{\zihao{5}注：开题报告用A4纸打印装订在毕业设计（论文）任务书后，学生可根据开题报告的长度加页。}

\end{document}
